\documentclass[11pt]{article}
\usepackage[margin=0.8in]{geometry}
\usepackage{booktabs,longtable,array}
\usepackage{float}
\usepackage{pdflscape}
\usepackage{siunitx}
\usepackage[T1]{fontenc}
\usepackage{lmodern}
\usepackage{hyperref}
\sisetup{detect-all}
\renewcommand{\thetable}{S\arabic{table}}
\renewcommand{\thefigure}{S\arabic{figure}}
\setlength{\LTleft}{0pt}
\setlength{\LTright}{0pt}

\title{Supporting Information for\\
BioInteract: Interpretable Drug--Target Interaction Prediction via Residue-Level Cross-Attention with Biological Prior Knowledge}
\author{ }
\date{ }

\begin{document}
\maketitle

\noindent This document reports only values recoverable from the released Davis
records, configuration files, saved result files, and source code included with
this revision. ``Classifier score'' denotes the sigmoid output for the binary
high-affinity class; it is not an experimental affinity. Model-native attention
and Grad-CAM values are attribution scores, not contact assignments.

\section*{Model configuration and input features}

\begin{table}[H]
\centering
\caption{Model architecture from \texttt{configs/default.yaml}.}
\begin{tabular}{ll}
\toprule
Component & Archived setting \\
\midrule
Drug encoder & GINE; 3 layers; hidden dimension 256; dropout 0.20 \\
Atom and bond input & 52 node features; 16 edge features \\
Protein language model & \texttt{esm2\_t30\_150M\_UR50D}; 640-dimensional residues \\
Protein projection & 256-dimensional projection; 32-dimensional domain-label embedding \\
Cross-attention & 8 heads; attention dimension 256; dropout 0.10 \\
Pooling & Gated \\
Prediction head & Hidden dimensions 256 and 128; dropout 0.30; classification task \\
\bottomrule
\end{tabular}
\end{table}

\begin{table}[H]
\centering
\caption{Protein-residue input channels and their availability.}
\begin{tabular}{p{0.26\linewidth}p{0.66\linewidth}}
\toprule
Channel & Archived implementation and limitation \\
\midrule
ESM-2 & Frozen 640-dimensional residue embeddings; proteins are truncated at 1,200 residues. \\
Physicochemical vector & Four values per residue: hydrophobicity, normalized molecular weight, charge, and polarity. \\
Domain-label channel & The configuration reserves a 32-dimensional embedding. The label builder uses \texttt{NONE} when no per-protein JSON annotation is available. \\
Annotation availability & No \texttt{data/domain\_annotations/*.json} files are present in the released package. Thus the reproducible Davis input uses the unknown label at every residue and is not curated Pfam/InterPro information. \\
\bottomrule
\end{tabular}
\end{table}

\begin{table}[H]
\centering
\caption{Training settings from \texttt{configs/default.yaml}.}
\begin{tabular}{ll}
\toprule
Setting & Value \\
\midrule
Dataset and task & Davis; binary high-affinity interaction classification \\
Random seed & 42 \\
Split proportions & 0.70 train / 0.10 validation / 0.20 test \\
Batch size & 64 \\
Optimizer & AdamW; learning rate $10^{-4}$; weight decay $10^{-5}$ \\
Schedule & Cosine; 5-epoch warmup \\
Maximum epochs and patience & 100 and 15 \\
Gradient accumulation and AMP & 1 and enabled \\
Protein maximum length & 1,200 residues \\
\bottomrule
\end{tabular}
\end{table}

\begin{table}[H]
\centering
\caption{Computational environment documented during final reproducibility audit.}
\begin{tabular}{p{0.30\linewidth}p{0.62\linewidth}}
\toprule
Item & Value \\
\midrule
Python & 3.10.14 \\
PyTorch & 2.4.0+cu121 \\
PyTorch Geometric & 2.6.1 \\
RDKit & 2025.3.6 \\
\texttt{fair-esm} & 2.0.0 \\
Biopython and RapidFuzz & 1.87 and 3.14.5 \\
CUDA availability during audit & Available \\
Audit GPU & NVIDIA GeForce RTX 4080 \\
Strict-split run record & \texttt{device: cuda} in \texttt{strict\_split\_metrics.json} \\
Checkpoint files available & \texttt{best.pt}, \texttt{best\_random.pt}, \texttt{best\_cold\_target.pt}, \texttt{best\_cold\_drug.pt} \\
\bottomrule
\end{tabular}
\end{table}

\clearpage

\section*{Dataset, partitions, and metrics}

\begin{table}[H]
\centering
\caption{Canonical current-release Davis test-partition summaries. Values are
from the full-precision checkpoint reconstruction, frozen validation thresholds,
and the associated prediction files.}
\begin{tabular}{lrrrrrrr}
\toprule
Partition & Test pairs & Positives & AUROC & AUPRC & F1 & Precision & Recall \\
\midrule
Random & 6,011 & 276 & 0.904 & 0.560 & 0.565 & 0.602 & 0.533 \\
Target-ID-held-out & 5,984 & 250 & 0.930 & 0.524 & 0.534 & 0.517 & 0.552 \\
Drug-ID-held-out & 5,746 & 245 & 0.733 & 0.166 & 0.100 & 0.273 & 0.061 \\
\bottomrule
\end{tabular}
\end{table}

\begin{table}[H]
\centering
\caption{Bootstrap intervals for the canonical current-release entity-based partitions.
Each interval uses 600 ordinary pair-level bootstrap replicates sampled with
replacement at the test-pair level, NumPy seed 42, and the 2.5th and 97.5th
percentiles of valid replicate metrics. Replicates containing only one class
are discarded. The main figure shows point estimates only.}
\begin{tabular}{lcc}
\toprule
Partition & AUROC 95\% interval & AUPRC 95\% interval \\
\midrule
Random & [0.881, 0.924] & [0.498, 0.624] \\
Target-ID-held-out & [0.911, 0.948] & [0.454, 0.584] \\
Drug-ID-held-out & [0.703, 0.763] & [0.129, 0.211] \\
\bottomrule
\end{tabular}
\end{table}

\begin{table}[H]
\centering
\caption{Exact-sequence-grouped strict within-Davis evaluations from \texttt{strict\_split\_metrics.json}. Threshold-dependent quantities use the corresponding validation-set threshold.}
\scriptsize
\textbf{Pair accounting}\par\smallskip
\begin{tabular}{lrrrrr}
\toprule
Protocol & Train pairs & Validation pairs & Test pairs & Test positives & Excluded cross-block pairs \\
\midrule
Exact-sequence-grouped cold-target & 21,080 & 2,992 & 5,984 & 208 & 0 \\
Exact-sequence-grouped cold-both & 15,190 & 264 & 1,144 & 38 & 13,458 \\
\bottomrule
\end{tabular}

\medskip
\textbf{Test metrics}\par\smallskip
\begin{tabular}{lrrrrrr}
\toprule
Protocol & AUROC & AUPRC & Threshold & F1 & Precision & Recall \\
\midrule
Exact-sequence-grouped cold-target & 0.873 & 0.383 & 0.838 & 0.468 & 0.517 & 0.428 \\
Exact-sequence-grouped cold-both & 0.552 & 0.038 & 0.936 & 0.000 & 0.000 & 0.000 \\
\bottomrule
\end{tabular}

\medskip
\raggedright For the cold-both protocol, drug identifiers and exact
target-sequence groups are partitioned independently at 70/10/20. Only pairs in
the corresponding train--train, validation--validation, and test--test blocks
are retained; the remaining cross-block pairs are excluded. The cold-both
validation block contains 10 positive pairs.
\end{table}

\begin{landscape}
\begin{table}[H]
\centering
\caption{Source and protocol context for baseline values reproduced from the
comparison table in the originally submitted manuscript. The table identifies
the numerical source, reported evaluation setting, protein-representation
context, and the appropriate descriptive interpretation of the comparison.}
\footnotesize
\setlength{\tabcolsep}{2.8pt}
\begin{tabular}{p{0.13\linewidth}p{0.19\linewidth}p{0.16\linewidth}p{0.16\linewidth}p{0.16\linewidth}p{0.15\linewidth}}
\toprule
Method & Numerical source & Original comparison protocol & Binary criterion & Protein representation & Interpretation in this revision \\
\midrule
DeepDTA & Original submitted comparison table & Random, Target-ID-held-out, and Drug-ID-held-out & \SI{30}{nM} Davis affinity criterion & Baseline-native representation; no matched frozen ESM-2 input & Descriptive contextual comparison; not a matched-pretraining or common-code benchmark \\
GraphDTA & Original submitted comparison table & Random, Target-ID-held-out, and Drug-ID-held-out & \SI{30}{nM} Davis affinity criterion & Baseline-native representation; no matched frozen ESM-2 input & Descriptive contextual comparison; not a matched-pretraining or common-code benchmark \\
AttentionDTA & Original submitted comparison table & Random, Target-ID-held-out, and Drug-ID-held-out & \SI{30}{nM} Davis affinity criterion & Baseline-native representation; no matched frozen ESM-2 input & Descriptive contextual comparison; not a matched-pretraining or common-code benchmark \\
MolTrans & Original submitted comparison table & Random, Target-ID-held-out, and Drug-ID-held-out & \SI{30}{nM} Davis affinity criterion & Baseline-native representation; no matched frozen ESM-2 input & Descriptive contextual comparison; not a matched-pretraining or common-code benchmark \\
TransformerCPI & Original submitted comparison table & Random, Target-ID-held-out, and Drug-ID-held-out & \SI{30}{nM} Davis affinity criterion & Baseline-native representation; no matched frozen ESM-2 input & Descriptive contextual comparison; not a matched-pretraining or common-code benchmark \\
DrugBAN & Original submitted comparison table & Random, Target-ID-held-out, and Drug-ID-held-out & \SI{30}{nM} Davis affinity criterion & Baseline-native representation; no matched frozen ESM-2 input & Descriptive contextual comparison; not a matched-pretraining or common-code benchmark \\
BioInteract & Archived prediction CSVs and result JSONs & Random, Target-ID-held-out, and Drug-ID-held-out & \SI{30}{nM} Davis affinity criterion & Frozen ESM-2 150M residue embeddings & Traceable within BioInteract; not a cross-method matched-encoder comparison \\
\bottomrule
\end{tabular}
\end{table}
\end{landscape}

\section*{Reproducibility details for revision analyses}

\noindent\textbf{Identifier-partition similarity audits.}
Drug similarities use RDKit~2025.3.6 Morgan fingerprints (radius 2, 1,024 bits,
\texttt{includeChirality=False}) and Tanimoto similarity. Protein similarities
for the original Target-ID-held-out audit use RapidFuzz~3.14.5
\texttt{fuzz.ratio}/100. For sequences $x,y$, this is the normalised Indel
similarity $1-d_{\mathrm{Indel}}/(|x|+|y|)$, where substitutions are represented
as one deletion plus one insertion.

\medskip
\noindent\textbf{Strict-test global alignment.}
Biopython~1.87 \texttt{Align.PairwiseAligner} is used in global mode with match
$=1$, mismatch $=0$, gap-opening $=-1$, and gap-extension $=-0.5$. Full-length
identity is the number of exact-residue matches divided by alignment length; the
maximum across all training targets is retained for each held-out target.

\medskip
\noindent\textbf{Bootstrap intervals.}
The saved 95\% intervals use ordinary, unstratified pair-level resampling with
replacement, 600 requested replicates, and NumPy random seed 42. Each replicate
contains the same number of test pairs as the corresponding test set; one-class
replicates are skipped. Bounds are the 2.5th and 97.5th percentiles of the valid
AUROC or AUPRC replicate distribution.

\medskip
\noindent\textbf{Grad-CAM.}
Hooks are attached to \texttt{drug\_encoder.layers.2}, the third and final GINE
message-passing layer. Gradients are taken from the pre-sigmoid positive-class
logit, pooled across graph nodes to obtain channel weights, rectified, and then
normalised separately within each molecular graph.

\medskip
\noindent\textbf{Attention sparsity.}
For pair $q$, residue scores $s_{qj}=\sum_i M_{qij}$ are normalised as
$\widetilde{s}_{qj}=s_{qj}/\max_k s_{qk}$. The reported fraction is
$|\{(q,j):\widetilde{s}_{qj}<0.1\}|/\sum_q L_q$. It uses all 1,506 positive Davis
pairs with the archived \texttt{checkpoints/best.pt}, not a single test set.
Every pair is normalised before the 1,301,380 valid residue scores are
concatenated; thus original training, validation, and test memberships are all
included and longer targets contribute more residue-level observations.

\begin{center}
\small
\textbf{Exact-sequence-grouped target diagnostic controls.} All values use the 5,984-pair strict target test set with 208 positives.\\[3pt]
\begin{tabular}{lrr}
\toprule
Control & AUROC & AUPRC \\
\midrule
BioInteract & 0.873 & 0.383 \\
Drug-only Morgan logistic regression & 0.764 & 0.164 \\
Drug-promiscuity prior & 0.764 & 0.163 \\
Target-only constant prior & 0.500 & 0.035 \\
Label-shuffled drug-only control & 0.407 & 0.029 \\
\bottomrule
\end{tabular}
\end{center}

\clearpage

\section*{Feature and input audits}

\begin{table}[H]
\centering
\caption{Atom and bond feature specification from \texttt{src/data/mol\_graph.py}.}
\begin{tabular}{llr}
\toprule
Level & Feature family & Dimensions \\
\midrule
Atom & Atom type one-hot (20 choices including unknown) & 20 \\
Atom & Hybridization, formal charge, and hydrogen-count one-hot vectors & 18 \\
Atom & Degree, aromaticity, and ring membership & 3 \\
Atom & Ring-size membership for sizes 3--8 & 6 \\
Atom & Donor/acceptor, Crippen logP and MR, Gasteiger charge & 5 \\
Atom total &  & 52 \\
Bond & Bond type, conjugation, ring membership, stereo, and direction & 16 \\
\bottomrule
\end{tabular}
\end{table}

\begin{longtable}{@{}l@{\quad}l@{\quad}r@{\quad}c@{\quad}r@{\quad}r@{}}
\caption{Davis nominal ABL1 input audit. Every listed record has the same complete SHA-256 digest {\scriptsize\texttt{cd6a78b82fede215cdb80a797f228098a41178fdde0753728f3781e6221fb1ee}}; values in the last columns are the archived D0010 and D0017 affinities in nM. The nominal labels do not alter the model sequence input.}\\
\toprule
Target ID & Nominal name & Length & Matches ABL1 & D0010 & D0017 \\
\midrule
\endfirsthead
\toprule
Target ID & Nominal name & Length & Matches ABL1 & D0010 & D0017 \\
\midrule
\endhead
T0001 & ABL1(E255K) & 1167 & Yes & 0.047 & 0.047 \\
T0002 & ABL1(F317I) & 1167 & Yes & 0.63 & 0.10 \\
T0003 & ABL1(F317I)p & 1167 & Yes & 0.18 & 0.041 \\
T0004 & ABL1(F317L) & 1167 & Yes & 0.11 & 0.032 \\
T0005 & ABL1(F317L)p & 1167 & Yes & 0.029 & 0.019 \\
T0006 & ABL1(H396P) & 1167 & Yes & 0.062 & 0.025 \\
T0007 & ABL1(H396P)p & 1167 & Yes & 0.057 & 0.046 \\
T0008 & ABL1(M351T) & 1167 & Yes & 0.037 & 0.016 \\
T0009 & ABL1(Q252H) & 1167 & Yes & 0.086 & 0.037 \\
T0010 & ABL1(Q252H)p & 1167 & Yes & 0.039 & 0.064 \\
T0011 & ABL1(T315I) & 1167 & Yes & 21 & 890 \\
T0012 & ABL1(T315I)p & 1167 & Yes & 3.6 & 120 \\
T0013 & ABL1(Y253F) & 1167 & Yes & 0.036 & 0.058 \\
T0014 & ABL1 & 1167 & Yes & 0.12 & 0.029 \\
T0015 & ABL1p & 1167 & Yes & 0.057 & 0.046 \\
\bottomrule
\end{longtable}

\begin{table}[H]
\centering
\caption{Saved non-mutant EGFR--D0014 functional-group attribution values.}
\begin{tabular}{lr}
\toprule
Functional group & Attribution \\
\midrule
Carbonyl & 0.431 \\
Amide & 0.389 \\
Ether & 0.136 \\
Halogen & 0.134 \\
Amino & 0.090 \\
Aromatic ring & 0.000 \\
Heterocycle N & 0.000 \\
\bottomrule
\end{tabular}
\end{table}

\begin{table}[H]
\centering
\caption{Saved non-mutant BRAF--D0023 failure-mode summary. The archived affinity is positive at the study threshold, whereas the classifier score is low. The score and attributions are model outputs, not an explanation of binding.}
\begin{tabular}{lr}
\toprule
Quantity & Value \\
\midrule
High-affinity-class score & 0.099 \\
Archived Davis affinity (nM) & 0.19 \\
Aromatic-ring attribution & 0.059 \\
Hydroxyl attribution & 0.000 \\
Heterocycle-N attribution & 0.000 \\
\bottomrule
\end{tabular}
\end{table}

\begin{table}[H]
\centering
\caption{Trainable parameter count recovered from \texttt{interpretability\_report.json}.}
\begin{tabular}{lr}
\toprule
Quantity & Value \\
\midrule
Total trainable parameters & 2,442,083 \\
\bottomrule
\end{tabular}
\end{table}

\end{document}
